% Generated draft: requires mathematical and source review.
\documentclass[11pt]{article}
\usepackage[margin=25mm]{geometry}
\usepackage{fontspec}
\setmainfont{Latin Modern Roman}
\usepackage{amsmath,amssymb,mathtools,unicode-math}
\setmathfont{Latin Modern Math}
\usepackage{enumitem,needspace,xurl,hyperref}
\setlength{\parindent}{0pt}
\setlength{\parskip}{5pt}
\setlength{\emergencystretch}{3em}
\newcommand{\N}{\mathbb{N}}
\newcommand{\Z}{\mathbb{Z}}
\newcommand{\Q}{\mathbb{Q}}
\newcommand{\R}{\mathbb{R}}
\newcommand{\C}{\mathbb{C}}
\newcommand{\F}{\mathbb{F}}
\newcommand{\A}{\mathbb{A}}
\newcommand{\PP}{\mathbb{P}}
\newcommand{\E}{\mathbb{E}}
\newcommand{\eps}{\varepsilon}
\newcommand{\dd}{\,\mathrm{d}}
\begin{document}
\section*{UnsolvedMath ID: 701778}
\textbf{Source problem code:} catalogue entry
\subsection{Definitions and mathematical statement}
A group \(G\) is locally nilpotent if every finitely generated subgroup \(H\leq G\) is nilpotent. Explicitly, writing \(\gamma_1(H)=H\) and \(\gamma_{i+1}(H)=[\gamma_i(H),H]\), this means that for each such \(H\) there exists an integer \(c\geq 1\) with \(\gamma_{c+1}(H)=\{1\}\). Here \([A,B]\) denotes the subgroup generated by the commutators \(a^{-1}b^{-1}ab\), with \(a\in A\) and \(b\in B\).

For an automorphism \(\alpha\in\operatorname{Aut}(G)\), interpret simplicity of \(G\) as a group with automorphism to mean that \(G\neq\{1\}\) and every normal subgroup \(N\trianglelefteq G\) satisfying \(\alpha(N)=N\) is either \(\{1\}\) or \(G\). This is the invariant-normal-subgroup interpretation of the source terminology; the supplied record does not explicitly define that terminology.

\textbf{Question (E. I. Zelmanov).} Does there exist a locally nilpotent group \(G\) together with an automorphism \(\alpha\in\operatorname{Aut}(G)\) of infinite order such that \(G\) is simple as a group with the distinguished automorphism \(\alpha\)? Equivalently, does there exist such a pair satisfying both
\[
\alpha^n\neq\operatorname{id}_G\qquad\text{for every integer }n\geq 1
\]
and
\[
\forall N\trianglelefteq G,\qquad
\alpha(N)=N\ \Longrightarrow\ N\in\bigl\{\{1\},G\bigr\}?
\]
The source's parenthetical qualification ``of infinite order'' is retained as a requirement on \(\alpha\).
\subsection{Short English statement}
Can a locally nilpotent group admit an automorphism of infinite order for which the only invariant normal subgroups are the trivial subgroup and the whole group? This uses the invariant-normal-subgroup interpretation of simplicity as a group with automorphism.
\subsection{Sources}
\begin{itemize}
\item Supplied local UnsolvedMath registry record: permanent ID \(701778\), exact source problem code \texttt{catalogue entry}, title \emph{Kourovka Notebook Problem 20.120}; upstream revision \texttt{d371ee917ac9647fbb06cd6fc529deec6f3dda82}.
\item The supplied attribution is to E. I. Zelmanov, proposed in 2022. The supplied background identifies \emph{The Kourovka Notebook (Unsolved Problems in Group Theory)}, 20th issue (2022), original problem 20.120, and the 21st issue edition (2026), \texttt{arXiv:1401.0300v48}. The supplied collection description names E. I. Khukhro and V. D. Mazurov as editors.
\item The registry labels the problem \emph{open} and records an update timestamp of \texttt{2026-10-16T14:00:00Z}. This is supplied registry information, not an independently verified assessment of the literature. No source was independently consulted for this draft.
\end{itemize}
\end{document}
